\subsection{Lorenz-96}
\label{sec:results-l96}

We organize the Lorenz-96 results as follows: first the integrator and Lyapunov
diagnostics that establish the chaotic time scale of the testbed; then the training of
the learned inverse observation operator $\Hinv$; then the choice of cold-start
initialization; then our central single-window result, the tuned-$\sigb$ full 4D-Var
sweep; and finally the sliding-window (cycling) experiments.

\subsubsection{Integrator and Lyapunov diagnostics}

Our PyTorch port integrates Lorenz-96 (grid size $40$,
forcing $F=8$) with a fixed-step RK4 scheme, ten inner
substeps of $\dtin = 0.01$ per outer step $\dtout = 0.1$, whereas the reference paper
uses an adaptive dopri5 solver. To check
that this choice does not contaminate the DA results, we propagate
the same warmed-up attractor state $200$ outer steps ($20$ MTU) with
both integrators and measure the $L^2$ state divergence
$\lVert \xstate_{\mathrm{RK4}}(t) - \xstate_{\mathrm{adapt}}(t)\rVert_2$
(Figure~\ref{fig:l96-integrator}). Because the global RK4 truncation error is
$O(\dtin^4) = O(10^{-8})$, well below the adaptive tolerances, the two integrators
agree to near machine precision while $t \ll T_L$ and only diverge once the tiny
per-step difference is amplified at the chaotic rate. The divergence first reaches the
order-one (state-norm) scale at $t^\ast = 6.30$ MTU for a single trajectory; across five
warmed-up states the median is $t^\ast = 7.70$ MTU (range $[6.50, 8.30]$), and the
divergence saturates near $28.5$ by $t = 19.9$ MTU once the trajectories fully
decorrelate. A direct two-trajectory separation fit (initial separation
$\epsilon_0 = 10^{-6}$, fit window $t \in [0.5, 5.0]$ MTU; Figure~\ref{fig:l96-lyapunov})
gives a leading growth rate $\lambda = 2.263\,\mathrm{MTU}^{-1}$ and Lyapunov time
$T_L = 1/\lambda = 0.442$ MTU. This is somewhat shorter than the paper's expected
$\sim\!0.6$ MTU; we report it honestly as a known consequence of the fixed-step
integrator being slightly too coarse to resolve the fastest unstable directions. The
practical upshot is that $t^\ast$ is the operational ceiling on assimilation-window
length over which the two simulators represent the same physics: our single-window
experiments use $T = 10$ outer steps ($=1.0$ MTU), comfortably inside $t^\ast$, so the
integrator choice does not affect the DA conclusions below.

\begin{figure}[!ht]
  \centering
  \includegraphics[width=0.58\linewidth]{L96_Integrator__cell08__out00.png}
  \caption{Fixed-step RK4 versus adaptive dopri5 integrator divergence
  on Lorenz-96. Vertical axis (log scale): the $L^2$ state difference
  $\lVert \xstate_{\mathrm{RK4}}(t)-\xstate_{\mathrm{adapt}}(t)\rVert_2$; horizontal axis:
  model time $t$ in MTU. The solid blue curve is the median over five warmed-up attractor
  states, the shaded band the min/max envelope, and the dashed grey line marks the
  order-one threshold. The two integrators agree to near machine precision while
  $t \ll T_L$, then diverge at the chaotic rate, crossing the order-one scale at a median
  $t^\ast = 7.70$ MTU and saturating near $28.5$. Our single-window experiments use
  $T = 1.0$ MTU, comfortably inside $t^\ast$, so the integrator choice does not affect the
  DA results.}
  \label{fig:l96-integrator}
\end{figure}

\begin{figure}[!ht]
  \centering
  \includegraphics[width=0.58\linewidth]{L96_Integrator__cell10__out00.png}
  \caption{Direct Lyapunov-time estimate for Lorenz-96. Two trajectories are launched
  $\epsilon_0 = 10^{-6}$ apart and the $L^2$ separation $\lVert\delta\xstate(t)\rVert_2$
  (orange, log scale) is plotted against model time $t$; the dashed black line is an
  exponential fit over the linear-growth window $t \in [0.5, 5.0]$ MTU and the dotted grey
  line marks the saturation scale. The fit yields growth rate
  $\lambda = 2.263\,\mathrm{MTU}^{-1}$ and Lyapunov time $T_L = 1/\lambda = 0.442$ MTU,
  after which the separation saturates near the attractor diameter. This is somewhat
  shorter than the paper's expected $\sim\!0.6$ MTU, a known consequence of the fixed-step
  RK4 integrator being slightly too coarse to resolve the fastest unstable directions.}
  \label{fig:l96-lyapunov}
\end{figure}

\subsubsection{Inverse-operator training}

The learned inverse observation operator $\Hinv$ maps the observation sequence $Y$ (only
$10$ of $40$ grid points, since $\Hop$ subsamples every fourth point) back to the full
physical-state sequence $X$. We train two CNN variants on
the same $32{,}000$-trajectory dataset with identical Adam settings ($500$ epochs, batch
size $256$, learning rate $10^{-3}$). The original PyTorch port (GELU
activations; flat width-$32$ residual stack; bilinear upsampling) reaches a
final training MSE of $0.3714$, while the paper-faithful inverter
(SiLU activations; channel pyramid
$128\!\to\!64\!\to\!32\!\to\!16$; BN after every convolution;
cubic-periodic upsampling; no skip connections) reaches $0.4924$
(Figure~\ref{fig:l96-invtrain}). On a held-out $100$-trajectory set (seed $100$), the
port attains test $L_1 = 0.3838$ with $\xstate_0$ RMSE $0.9973 \pm 0.3261$, and the
paper-faithful inverter attains test $L_1 = 0.4470$ with $\xstate_0$ RMSE
$1.1635 \pm 0.4036$. The paper/port $L_1$ ratio of $1.165$ lies comfortably within a
$30\%$ acceptance band, so the paper-faithful architecture reproduces the port's
reconstruction quality despite the architectural differences. Retraining the
paper-faithful inverter at higher noise gives final MSE $0.6321$, $1.3338$, and $2.2201$
at $\sigobs = 0.1, 0.5, 1.0$ respectively; these noise-matched checkpoints supply the
inverter used in the hybrid DA runs below.

\begin{figure}[!ht]
  \centering
  \includegraphics[width=0.58\linewidth]{L96_PaperFaithful__cell10__out00.png}
  \caption{Training-loss curves for the two inverse observation operators at
  $\sigobs = 0$. Both networks are trained on the same $32{,}000$-trajectory dataset with
  identical Adam settings ($500$ epochs); vertical axis: training MSE (log scale);
  horizontal axis: epoch. The blue curve is the original PyTorch port (GELU activations,
  flat width-$32$ residual stack, bilinear upsampling), converging to final MSE $0.3714$;
  the orange curve is the paper-faithful inverter (SiLU activations,
  $128\!\to\!64\!\to\!32\!\to\!16$ channels, BN after every convolution, cubic-periodic
  upsampling, no skip connections), converging to $0.4924$. The two curves track closely,
  and despite the architectural differences the paper-faithful inverter reproduces the
  port's reconstruction quality (held-out test $L_1$ ratio $1.165$, within the $30\%$
  acceptance band).}
  \label{fig:l96-invtrain}
\end{figure}

\subsubsection{Baseline initialization}

A 4D-Var solve needs a cold-start initial guess for the optimizer. We compare the
original ``repeat-interleave'' baseline (copy each $t = 0$ observation across its block of
four unobserved points, producing step discontinuities at observation boundaries) against
the paper's climatological-mean baseline (fill unobserved positions with the attractor
mean $\bar X \approx 2.34$, then overwrite observed positions with $\yobs_0$). At
$\sigobs = 0$ ($T = 10$, $N = 100$), the raw initializers give direct $\xstate_0$ RMSE
$4.7828 \pm 0.5006$ (repeat) versus $3.0892 \pm 0.2539$ (climatological mean), a ratio of
$0.646$, i.e.\ the climatological-mean prior is roughly $35\%$ closer to truth. The
spatial-correlation matrix $\bm{C}$ used for decorrelation preconditioning has condition
number $8.72$. Figure~\ref{fig:l96-baselinebar} shows analysis RMSE for the four
(initializer $\times$ optimization-mode) combinations under untuned full 4D-Var
($\sigb = 1.0$ fixed): both a better cold-start prior and hybrid optimization lower the
analysis error before $\sigb$ is even tuned, with the climatological-mean baseline plus
hybrid optimization lowest at every noise level. The space-time (Hovm\"oller) comparison
at $\sigobs = 0.5$ (Figure~\ref{fig:l96-hovmoller}) makes the qualitative difference
concrete: observation-space 4D-Var from the repeat-interleave baseline carries visible
artifacts and a distorted wave pattern out of the assimilation window, whereas the
climatological-mean baseline reproduces the truth's traveling-wave structure far more
faithfully into the forecast.

\begin{figure}[!ht]
  \centering
  \includegraphics[width=0.82\linewidth]{L96_PaperFaithful__cell19__out00.png}
  \caption{Analysis RMSE comparison of the two baseline initializations under untuned full
  4D-Var ($\sigb = 1.0$ fixed). Grouped bars show post-assimilation state RMSE on
  $\xstate_0$ ($\pm$ one standard deviation over $N = 100$ trajectories) for the four
  combinations of initializer (repeat-interleave vs.\ climatological-mean) and
  optimization mode (observation-space vs.\ hybrid), across noise levels
  $\sigobs \in \{0.1, 0.5, 1.0\}$. The repeat-interleave baseline with observation-only
  optimization (grey) is worst at every noise level; the climatological-mean baseline with
  hybrid optimization is lowest (most clearly at $\sigobs = 0.1$). A better cold-start
  prior and hybrid optimization each reduce analysis error even before $\sigb$ is tuned,
  motivating the per-cell $\sigb$ sweep that follows.}
  \label{fig:l96-baselinebar}
\end{figure}

\begin{figure}[!ht]
  \centering
  \includegraphics[width=0.85\linewidth]{L96_PaperFaithful__cell20__out00.png}
  \caption{Hovm\"oller comparison of the two baselines at $\sigobs = 0.5$. Each panel
  shows the Lorenz-96 state $x_k$ (color, red--blue diverging scale) with grid index $k$
  on the vertical axis and time step on the horizontal axis: ground-truth trajectory
  (left), observation-space 4D-Var initialized from the repeat-interleave baseline
  (center), and from the climatological-mean baseline (right). The dashed vertical line
  marks the end of the $10$-step assimilation window; everything to its right is free
  forecast. The repeat baseline carries visible artifacts and a distorted wave pattern out
  of the window, whereas the climatological-mean baseline reproduces the truth's
  traveling-wave structure more faithfully into the forecast under moderate observation
  noise.}
  \label{fig:l96-hovmoller}
\end{figure}

\subsubsection{Tuned-\texorpdfstring{$\sigb$}{sigma\_b} full 4D-Var: the key result}

Our central single-window experiment sweeps the background-error standard deviation
$\sigb$ over $\{0.1, 0.3, 1.0, 3.0, 10.0\}$ for every (initializer $\times$
optimization-mode $\times$ noise) cell of the full 4D-Var cost
$J = \Jb + \Jo$, with $\Bcov = \sigb^2 \bm{C}$ and $\Rcov = \sigobs^2 \bm{I}$, and records
the lowest-RMSE choice ($N = 16$, $T = 10$). Table~\ref{tab:l96-tuned} reports the best
analysis RMSE per cell with the selected $\sigb$ in parentheses, and
Figure~\ref{fig:l96-keyresult} plots the same data with the optimal $\sigb$ annotated.

\begin{table}[H]
  \centering
  \caption{Best analysis RMSE on $\xstate_0$ for tuned-$\sigb$ full 4D-Var on Lorenz-96
  ($N = 16$, $T = 10$), with the selected $\sigb \in \{0.1, 0.3, 1.0, 3.0, 10.0\}$ in
  parentheses. ``paper'' denotes the climatological-mean baseline and ``invobs'' the
  learned-inverse initialization; ``obs'' is observation-space optimization and ``hybrid''
  the physics-space-then-observation-space schedule. Bold marks the lowest RMSE in each
  noise column.}
  \label{tab:l96-tuned}
  \begin{tabular}{lccc}
    \toprule
    Method & $\sigobs = 0.1$ & $\sigobs = 0.5$ & $\sigobs = 1.0$ \\
    \midrule
    paper $+$ obs    & $2.466\ (0.1)$           & $2.393\ (0.3)$           & $2.708\ (0.3)$ \\
    paper $+$ hybrid & $0.621\ (3.0)$           & $1.223\ (1.0)$           & $1.846\ (1.0)$ \\
    invobs $+$ obs   & $0.883\ (0.3)$           & $1.550\ (0.3)$           & $1.867\ (1.0)$ \\
    invobs $+$ hybrid& $\mathbf{0.620}\ (10.0)$ & $\mathbf{1.168}\ (1.0)$  & $\mathbf{1.676}\ (1.0)$ \\
    \bottomrule
  \end{tabular}
\end{table}

Three findings stand out. First, the positive headline: invobs initialization combined
with hybrid (physics-space then observation-space) optimization attains the lowest
analysis RMSE at every noise level ($0.620$, $1.168$, $1.676$ for
$\sigobs = 0.1, 0.5, 1.0$), confirming that the paper's central idea survives in our full
4D-Var setting with observation noise (at $\sigobs = 0.1$ this is effectively a tie with
paper$+$hybrid, $0.620$ vs $0.621$). Second, hybrid optimization dramatically beats
observation-only optimization for the climatological-mean baseline that has no learned
lift: at $\sigobs = 0.5$, paper$+$obs ($2.393$) is roughly halved to paper$+$hybrid
($1.223$), because observation-only 4D-Var from a poor prior falls into bad local minima
when only $10$ of $40$ points are observed, and the physics-space warm start escapes them.
Third, an honest negative result: the expected ordering---that invobs$+$hybrid should
prefer a \emph{smaller} optimal $\sigb$ than paper$+$obs---does not hold. invobs$+$hybrid
selects $\sigb = 10.0/1.0/1.0$ while paper$+$obs selects $\sigb = 0.1/0.3/0.3$, the
opposite direction, failing the acceptance check at all three noise levels. We also note a
secondary caveat: tuning $\sigb$ helped only marginally over the untuned $\sigb = 1$ in
most cells (gains typically below $0.2$ RMSE, and exactly zero where the best $\sigb$ was
already $1.0$), and the small ensemble ($N = 16$) makes the error bars wide.

\begin{figure}[!ht]
  \centering
  \includegraphics[width=0.88\linewidth]{L96_PaperFaithful__cell28__out00.png}
  \caption{KEY RESULT---tuned-$\sigb$ full 4D-Var on Lorenz-96. One panel per
  observation-noise level ($\sigobs = 0.1, 0.5, 1.0$, left to right); within each panel the
  four bars are the (initializer $\times$ optimization-mode) combinations paper$+$obs,
  paper$+$hybrid, invobs$+$obs, invobs$+$hybrid, showing the lowest analysis RMSE on
  $\xstate_0$ ($\pm$ one standard deviation, $N = 16$) after sweeping
  $\sigb \in \{0.1, 0.3, 1.0, 3.0, 10.0\}$; the optimal $\sigb$ is annotated above each
  bar. The learned-inverse initialization with hybrid optimization (rightmost bar) attains
  the lowest RMSE at every noise level ($0.620$, $1.168$, $1.676$), and hybrid optimization
  dramatically improves the climatological-mean baseline over observation-only (paper$+$obs
  $2.393 \to$ paper$+$hybrid $1.223$ at $\sigobs = 0.5$). The annotated optimal $\sigb$
  values also show the honest negative result discussed in the text: invobs$+$hybrid
  selects $\sigb = 10.0/1.0/1.0$ versus paper$+$obs $0.1/0.3/0.3$, the opposite of the
  expected ordering.}
  \label{fig:l96-keyresult}
\end{figure}

\subsubsection{Sliding-window (cycling) 4D-Var}

We next extend the single fixed window to an operational cycling scheme. All Lorenz-96
sliding-window experiments use grid size $40$, $\Hop$ observing every fourth point,
per-cycle window $\texttt{WINDOW\_T} = 8$ steps, a $50$-step forecast horizon evaluated
after the last window, and an ensemble of $N = 50$ trajectories at noise levels
$\sigobs \in \{0.0, 0.5\}$. Single-window (cold-start) runs use background-error standard
deviation $\sigb = 1.0$, while cycling runs use a tighter $\sigb = 0.3$ because the cycled
background---the previous analysis propagated forward to the new window start---is
trusted more than a cold start. Optimization is L-BFGS in
decorrelated coordinates $\bm{z} = \bm{C}^{-1/2}\xstate$ ($\bm{C}$ condition number
$8.72$); hybrid mode runs $100$ physics-space iterations then $400$ observation-space
iterations. The roles of ``L-BFGS starting point'' and ``$\Jb$ background'' are
decoupled: every cycle restarts the optimizer from a fresh $\Hinv$ estimate of the new
window, while for cycles $\ge 1$ the $\Jb$ background is the propagated previous analysis
(cycle $0$ is a cold start whose background equals its starting point).

\paragraph{Experiment A (single 8-step window).}
Figure~\ref{fig:l96-swA} assimilates one $8$-step window at $\sigb = 1.0$ for the four
(initialization $\times$ optimization) combinations. All curves drop to low RMSE inside
the window and then grow toward the Lorenz-96 climatological saturation
($\rmse \approx 5$) over the chaotic forecast, consistent with the short Lyapunov time
($T_L \approx 0.44$ MTU). The analysis-RMSE bars at $t = 0$ are the clean discriminator:
the three optimization-capable methods reach analysis RMSE $\approx 1$ at $\sigobs = 0$
and $\approx 1.6$--$1.9$ at $\sigobs = 0.5$, while baseline$+$obs-only stays near $5$ at
both noise levels---the same observation-space local-minimum failure documented above,
which hybrid (physics-space) warm-starting rescues. By lowest mean forecast RMSE the best
single-window method is baseline$+$hybrid at $\sigobs = 0$ and invobs$+$hybrid at
$\sigobs = 0.5$; these selections feed the references in the experiments below. The
$\sigobs = 0$ margin among the three good methods is within ensemble scatter, so that
``best'' label is a narrow tie-break.

\begin{figure}[!ht]
  \centering
  \includegraphics[width=0.85\linewidth]{L96_SlidingWindow__cell27__out00.png}
  \caption{Experiment A: single $8$-step assimilation window on Lorenz-96 ($N = 50$,
  $\sigb = 1.0$), comparing four initialization $\times$ optimization combinations.
  \textbf{Top row:} ensemble-mean spatial $\rmse$ versus time step for $\sigobs = 0.0$
  (left) and $\sigobs = 0.5$ (right); the shaded band marks the assimilation window (steps
  $0$--$8$), the dashed line the window end, and coloured envelopes are $\pm 1$ ensemble
  standard deviation. All curves fall to low error inside the window and then grow toward
  the climatological saturation ($\rmse \approx 5$) over the $50$-step forecast, except
  baseline$+$obs-only, which never leaves the climatological level. \textbf{Bottom:}
  grouped bar chart of analysis $\rmse$ at the window start ($t = 0$): the three
  optimization-capable methods reach $\approx 1$ at $\sigobs = 0$ and $\approx 1.6$--$1.9$
  at $\sigobs = 0.5$, while baseline$+$obs-only stays near $5$. Hybrid warm-starting is
  essential to rescue the climatological-mean baseline from a poor observation-space local
  minimum.}
  \label{fig:l96-swA}
\end{figure}

\paragraph{Experiment B1 (stride sweep).}
For each observation history $T_\mathrm{obs} \in \{24, 48\}$ and
$\sigobs \in \{0, 0.5\}$ we run cycling with strides $\{2, 4, 8\}$ (cycle counts
$\{9, 5, 3\}$ for $T_\mathrm{obs} = 24$ and $\{21, 11, 6\}$ for $T_\mathrm{obs} = 48$), all
with invobs init each cycle, hybrid optimization, propagated background, and
$\sigb = 0.3$. Figure~\ref{fig:l96-swB1} shows the three stride curves are nearly
indistinguishable within ensemble scatter and all relax to the same forecast saturation;
the differences are confined to the window and shortest lead times. By mean forecast RMSE
the selected best strides are $8, 2, 8, 8$ for
$(T_\mathrm{obs}, \sigobs) = (24, 0), (24, 0.5), (48, 0), (48, 0.5)$---narrow tie-breaks
among overlapping curves rather than strong preferences. The dashed single-window
reference sits on top of the cycling curves over most of the forecast, foreshadowing that
cycling's advantage is modest and concentrated near the analysis time.

\begin{figure}[!ht]
  \centering
  \includegraphics[width=0.82\linewidth]{L96_SlidingWindow__cell30__out00.png}
  \caption{Experiment B1: sliding-window (cycling) 4D-Var stride sweep on Lorenz-96
  ($N = 50$, invobs init each cycle, hybrid optimization, propagated-analysis background,
  cycling $\sigb = 0.3$). Each panel plots last-cycle ensemble-mean spatial $\rmse$ versus
  absolute time step for strides $2, 4, 8$ (coloured) with $\pm 1$-std envelopes; rows are
  $T_\mathrm{obs} = 24$ (top) and $48$ (bottom), columns are $\sigobs = 0.0$ (left) and
  $0.5$ (right). The shaded band is the last assimilation window
  $[T_\mathrm{obs} - 8,\, T_\mathrm{obs}]$ and the dashed line is the window end; the dashed
  black curve is the best single-window Experiment-A method (baseline$+$hybrid at
  $\sigobs = 0$, invobs$+$hybrid at $\sigobs = 0.5$), time-shifted so its window end
  coincides with the cycling runs. The three stride curves are nearly indistinguishable
  within ensemble scatter and all relax to climatological saturation; by mean forecast
  $\rmse$ the selected best strides are $8, 2, 8, 8$. Assimilation stride has little effect
  on Lorenz-96 cycling skill, so the best stride is a narrow tie-break.}
  \label{fig:l96-swB1}
\end{figure}

\paragraph{Experiment B2 (per-cycle error reduction).}
For each panel's best-stride run, Figure~\ref{fig:l96-swB2} plots window-start RMSE per
cycle for the $\Jb$ background (the propagated previous analysis) and the analysis, with
the single-window invobs$+$hybrid analysis RMSE as a horizontal reference. The intended
mechanism is visible: the analysis sits at or below the background within each cycle (the
assimilation reduces error), the background error generally falls as cycling accumulates
history, and after the first cycle the cycled analysis drops below the single-window
reference. The descent is cleanest at $\sigobs = 0.5$ (most clearly $T_\mathrm{obs} = 48$,
stride $8$, where both curves fall monotonically and the analysis ends well under the
reference). The $\sigobs = 0$ panels are noisier and non-monotone (e.g.\ a mid-sequence
bump for $T_\mathrm{obs} = 48$): with perfect observations a single $8$-step window already
over-determines the $10$ observed components, so extra cycles add little and $N = 50$
sampling noise dominates.

\begin{figure}[!ht]
  \centering
  \includegraphics[width=0.8\linewidth]{L96_SlidingWindow__cell31__out00.png}
  \caption{Experiment B2: per-cycle error reduction for the best-stride sliding-window run
  in each $(T_\mathrm{obs}, \sigobs)$ configuration ($N = 50$). Each panel plots
  ensemble-mean spatial $\rmse$ at the window-start time versus cycle index for the $\Jb$
  background (red dashed squares, the previous analysis propagated to the new window) and
  the analysis (blue circles), with the single-window invobs$+$hybrid analysis $\rmse$ as a
  dotted horizontal reference. Panels: $T_\mathrm{obs} = 24$ best stride $8$
  ($\sigobs = 0$, $3$ cycles) and $2$ ($\sigobs = 0.5$, $9$ cycles); $T_\mathrm{obs} = 48$
  best stride $8$ for both noise levels ($6$ cycles). The analysis lies at or below the
  background within each cycle, and after cycle $0$ the cycled analysis drops below the
  single-window reference. The descent is cleanest at $\sigobs = 0.5$; at $\sigobs = 0$ the
  curves are noisier and non-monotone because a single $8$-step window already
  over-determines the $10$ observed components, leaving $N = 50$ sampling noise dominant.
  Cycling reduces background error into a better analysis each cycle, most clearly when
  observations are noisy.}
  \label{fig:l96-swB2}
\end{figure}

\paragraph{Experiment C (the key cycling comparison).}
Figure~\ref{fig:l96-swC} is the central cycling comparison: SW-best (best-stride cycling
with full observation history) versus two single-window methods that see only the last
$8$ observations---invobs$+$hybrid and baseline$+$obs-only---across the
$(T_\mathrm{obs}, \sigobs)$ grid. baseline$+$obs-only is consistently the worst, stuck near
the climatological RMSE through the window (the observation-space local-minimum failure
again). SW-best's per-cycle analysis chain descends as history accumulates, so it enters
the last-window analysis time from a lower error than the cold single-window methods. The
honest, central finding is that whether cycling \emph{clearly} wins depends on
$(T_\mathrm{obs}, \sigobs)$: the advantage is most pronounced at $\sigobs = 0.5$ with
longer history $T_\mathrm{obs} = 48$ (more cycles to average down observation noise and
propagate information), and marginal at $\sigobs = 0$, where a single well-posed $8$-step
window already pins the observed state and the methods overlap within $N = 50$ scatter.
Once into the chaotic forecast all methods converge to the same saturation regardless:
cycling improves the analysis and short-lead forecast, not the Lyapunov-time-limited
long-range predictability. This directly answers the experiment's question---cycling with
more history beats a single long/last window mainly when observations are noisy and history
is long.

\begin{figure}[!ht]
  \centering
  \includegraphics[width=0.88\linewidth]{L96_SlidingWindow__cell34__out00.png}
  \caption{Experiment C (key comparison): sliding-window cycling versus single-window
  4D-Var on Lorenz-96 ($N = 50$). Each $(T_\mathrm{obs}, \sigobs)$ panel (rows
  $T_\mathrm{obs} = 24, 48$; columns $\sigobs = 0.0, 0.5$) plots ensemble-mean spatial
  $\rmse$ versus absolute time step with $\pm 1$-std envelopes for three methods: SW-best
  (best-stride cycling with full observation history), a dashed per-cycle analysis-$\rmse$
  chain from $t = 0$ that hands off to a solid continuous curve for the last cycle plus
  $50$-step forecast; invobs$+$hybrid on the last $8$ observations; and baseline$+$obs-only
  on the last $8$ observations. The shaded band is the last assimilation window
  $[T_\mathrm{obs} - 8,\, T_\mathrm{obs}]$ and the dashed line marks the window end.
  baseline$+$obs-only stays near the climatological $\rmse$ throughout the window; SW-best's
  chain descends as history accumulates, so it enters the last-window analysis time from a
  lower error than the cold single-window methods. The cycling advantage is most pronounced
  at $\sigobs = 0.5$ with long history $T_\mathrm{obs} = 48$, and marginal at $\sigobs = 0$.
  Deep in the chaotic forecast all methods converge to the same climatological saturation:
  cycling improves the analysis and short-lead forecast, not the long-range predictability.}
  \label{fig:l96-swC}
\end{figure}

\paragraph{Experiment D (Hovm\"oller diagnostics).}
Finally, Figure~\ref{fig:l96-swE} shows space-time (Hovm\"oller) diagnostics for one
ensemble member across the four $(T_\mathrm{obs}, \sigobs)$ configurations, comparing the
truth, the SW-best reconstruction and signed error, and the single-window invobs$+$hybrid
reconstruction and signed error. Both methods reproduce the truth's traveling
wave packets through and just past the window; the signed-error rows are near-white inside
and immediately after the window and develop structured stripes as the chaotic forecast
decorrelates. SW-best shows visibly cleaner near-window error than the cold single-window
invobs$+$hybrid at $\sigobs = 0.5$, illustrating the same noise-dependent cycling advantage
seen in Experiment C, while error growth far into the forecast is comparable across
methods. This is a single illustrative member, not a statistical summary.

\begin{figure}[!ht]
  \centering
  \includegraphics[width=0.88\linewidth]{L96_SlidingWindow__cell36__out00.png}
  \caption{Experiment D: space-time (Hovm\"oller) diagnostics for one ensemble member
  (sample $0$) of the Lorenz-96 sliding-window run. Columns are the four configurations
  $(T_\mathrm{obs}, \sigobs) \in \{24, 48\} \times \{0.0, 0.5\}$; the five rows are
  (1) truth, (2) SW-best reconstructed state, (3) SW-best signed error (prediction $-$
  truth), (4) invobs$+$hybrid (last-$8$) reconstructed state, and (5) invobs$+$hybrid
  signed error, with grid index $k$ on the vertical axis and time step on the horizontal
  axis. State rows use a diverging colormap symmetric about $\pm\max|\text{truth}|$; the two
  error rows share a scale set to the larger $\max|\text{error}|$ across both methods within
  each column (blue $=$ under-prediction, white $=$ near-exact, red $=$ over-prediction),
  making the two methods directly comparable. The dashed vertical line marks the window end
  $T_\mathrm{obs}$: left of it is the last assimilation window, right of it the $50$-step
  forecast. Both methods reproduce the truth's travelling wave packets through and just past
  the window; the signed-error rows are near-white inside the window and develop structured
  stripes as the chaotic forecast decorrelates, with SW-best visibly cleaner near the window
  at $\sigobs = 0.5$. This is a single illustrative member ($N = 50$ overall), not a
  statistical summary.}
  \label{fig:l96-swE}
\end{figure}
